% GENERATED FILE. Edit canonical lessons, then run npm run generate:books.
% canonical-source-hash: fnv1a64:ae010fff464bcb60
% canonical-lessons: AR-C244-zabadi, AR-C244-khawkh, AR-C244-kummathra, AR-C244-ananas, AR-C244-lawz

\chapter{Yoghurt, Peaches, Pears, Pineapple, Almonds}
\label{ch:ar-yoghurt-peaches-pears-pineapple-almonds}
\hlchaptermodality{244}

\begin{chapteropening}
\textbf{By the end of this chapter:} \emph{I can say yoghurt, peaches, pears, pineapple and almonds.}

Complete the last lesson of chapter 244: I can say yoghurt, peaches, pears, pineapple and almonds.
\end{chapteropening}

\section[\texorpdfstring{\ar{زبادي} (zabādī)}{زبادي (zabādī)}]{\ar{زبادي} (zabādī) — yoghurt}

\label{lesson:AR-C244-zabadi}



\begin{quote}
Before the new one: say the Arabic for second-hand, then the Arabic for fake.
\end{quote}

\subsection*{You'll want to know: \ar{زبادي}}
\textbf{\ar{زبادي}} — \emph{zabādī} — \textquotedblleft{}yoghurt\textquotedblright{}.

\begin{cousinweb}
\textbf{\ar{زبادي}} (\emph{zabādī}) is a modern word for yoghurt. It may be linked with the root \emph{z-b-d}, which gives \emph{zubd}, \textquotedblleft{}fresh butter\textquotedblright{}, and \emph{zubād}, \textquotedblleft{}thin milk\textquotedblright{} (Lane), but its formation is not settled.
\end{cousinweb}

\begin{grammarlens}[title={What you've built}]
One more word for this chapter.
\end{grammarlens}

\subsection*{Your turn}
\begin{itemize}
\raggedright
  \item \emph{Say it:} \emph{zabādī}
  \item \emph{Say it:} \emph{zabādī}, once more
  \item \emph{Say it:} say \emph{zabādī}
  \item \emph{Recall:} say the Arabic for second-hand, then the Arabic for fake, then say \emph{zabādī} again
\end{itemize}

\begin{tcolorbox}[breakable,colback=teal!4,colframe=teal!35!black,title={Before you move on}]
What does \ar{زبادي} mean? (\textquotedblleft{}Yoghurt\textquotedblright{}.) Say it once more.
\end{tcolorbox}
\section[\texorpdfstring{\ar{خوخ} (khawkh)}{خوخ (khawkh)}]{\ar{خوخ} (khawkh) — peaches}

\label{lesson:AR-C244-khawkh}



\begin{quote}
Before the new one: say the Arabic for fake, then the Arabic for yoghurt.
\end{quote}

\subsection*{You'll want to know: \ar{خوخ}}
\textbf{\ar{خوخ}} — \emph{khawkh} — \textquotedblleft{}peaches\textquotedblright{}.

\begin{cousinweb}
\textbf{\ar{خوخ}} (\emph{khawkh}) is listed by Lane under \emph{kh-w-kh} as a well-known fruit, the peach, with no foreign source given; EtymDB derives it from Aramaic \emph{ḥūḥā}.
\end{cousinweb}

\begin{grammarlens}[title={What you've built}]
One more word for this chapter.
\end{grammarlens}

\subsection*{Your turn}
\begin{itemize}
\raggedright
  \item \emph{Say it:} \emph{khawkh}
  \item \emph{Say it:} \emph{khawkh}, once more
  \item \emph{Say it:} say \emph{khawkh}
  \item \emph{Recall:} say the Arabic for fake, then the Arabic for yoghurt, then say \emph{khawkh} again
\end{itemize}

\begin{tcolorbox}[breakable,colback=teal!4,colframe=teal!35!black,title={Before you move on}]
What does \ar{خوخ} mean? (\textquotedblleft{}Peaches\textquotedblright{}.) Say it once more.
\end{tcolorbox}
\section[\texorpdfstring{\ar{كمثرى} (kummathrā)}{كمثرى (kummathrā)}]{\ar{كمثرى} (kummathrā) — pears}

\label{lesson:AR-C244-kummathra}



\begin{quote}
Before the new one: say the Arabic for yoghurt, then the Arabic for peaches.
\end{quote}

\subsection*{You'll want to know: \ar{كمثرى}}
\textbf{\ar{كمثرى}} — \emph{kummathrā} — \textquotedblleft{}pears\textquotedblright{}.

\begin{cousinweb}
\textbf{\ar{كمثرى}} (\emph{kummathrā}) is usually traced to Syriac \emph{kumathrā}, \textquotedblleft{}pear\textquotedblright{} (EtymDB); it has no Arabic root.
\end{cousinweb}

\begin{grammarlens}[title={What you've built}]
One more word for this chapter.
\end{grammarlens}

\subsection*{Your turn}
\begin{itemize}
\raggedright
  \item \emph{Say it:} \emph{kummathrā}
  \item \emph{Say it:} \emph{kummathrā}, once more
  \item \emph{Say it:} say \emph{kummathrā}
  \item \emph{Recall:} say the Arabic for yoghurt, then the Arabic for peaches, then say \emph{kummathrā} again
\end{itemize}

\begin{tcolorbox}[breakable,colback=teal!4,colframe=teal!35!black,title={Before you move on}]
What does \ar{كمثرى} mean? (\textquotedblleft{}Pears\textquotedblright{}.) Say it once more.
\end{tcolorbox}
\section[\texorpdfstring{\ar{أناناس} (anānās)}{أناناس (anānās)}]{\ar{أناناس} (anānās) — pineapple}

\label{lesson:AR-C244-ananas}



\begin{quote}
Before the new one: say the Arabic for peaches, then the Arabic for pears.
\end{quote}

\subsection*{You'll want to know: \ar{أناناس}}
\textbf{\ar{أناناس}} — \emph{anānās} — \textquotedblleft{}pineapple\textquotedblright{}.

\begin{cousinweb}
\textbf{\ar{أناناس}} (\emph{anānās}) is a modern loan from European \emph{ananas}, which goes back through Portuguese to Tupi \emph{nanas}, the pineapple; it has no Arabic root.
\end{cousinweb}

\begin{grammarlens}[title={What you've built}]
One more word for this chapter.
\end{grammarlens}

\subsection*{Your turn}
\begin{itemize}
\raggedright
  \item \emph{Say it:} \emph{anānās}
  \item \emph{Say it:} \emph{anānās}, once more
  \item \emph{Say it:} say \emph{anānās}
  \item \emph{Recall:} say the Arabic for peaches, then the Arabic for pears, then say \emph{anānās} again
\end{itemize}

\begin{tcolorbox}[breakable,colback=teal!4,colframe=teal!35!black,title={Before you move on}]
What does \ar{أناناس} mean? (\textquotedblleft{}Pineapple\textquotedblright{}.) Say it once more.
\end{tcolorbox}
\section[\texorpdfstring{\ar{لوز} (lawz)}{لوز (lawz)}]{\ar{لوز} (lawz) — almonds}

\label{lesson:AR-C244-lawz}



\begin{quote}
Before the new one: say the Arabic for pears, then the Arabic for pineapple.
\end{quote}

\subsection*{You'll want to know: \ar{لوز}}
\textbf{\ar{لوز}} — \emph{lawz} — \textquotedblleft{}almonds\textquotedblright{}.

\begin{cousinweb}
\textbf{\ar{لوز}} (\emph{lawz}) is built on the native Arabic root \emph{l-w-z}, \textquotedblleft{}the almond\textquotedblright{} (Lane). Lane's \emph{lawz} is the almond; Hebrew \emph{luz} is a cousin.
\end{cousinweb}

\begin{grammarlens}[title={What you've built}]
One more word for this chapter.
\end{grammarlens}

\subsection*{Your turn}
\begin{itemize}
\raggedright
  \item \emph{Say it:} \emph{lawz}
  \item \emph{Say it:} \emph{lawz}, once more
  \item \emph{Say it:} say \emph{lawz}
  \item \emph{Recall:} say the Arabic for pears, then the Arabic for pineapple, then say \emph{lawz} again
\end{itemize}

\begin{tcolorbox}[breakable,colback=teal!4,colframe=teal!35!black,title={Before you move on}]
What does \ar{لوز} mean? (\textquotedblleft{}Almonds\textquotedblright{}.) Say it once more.
\end{tcolorbox}
